%% kapa: exam=1 round=30 year=2019 period=2 form=A subject=appraisal_law subject_ko=감정평가관계법규 seq=32 number=32
%% source: Q-Net 감정평가사 기출문제 (cst00309) artlSeq=5205205 file=extracted/1008/5205205/2019년 제30회 감정평가사 1차시험 문제지/2019년 제30회 감정평가사 1차시험 문제지/2019년도 제30회 감정평가사 1차 2교시 A형.hwp
%% conversion: kordoc 4.24.51 -> tools/gen1.py (auto); review=proofread-sonnet; flags=none
\begin{question}{32}
부동산등기법령상 ‘이의’에 관한 설명으로 옳지 않은 것은?
\begin{choices}
  \item 등기관의 결정 또는 처분에 이의가 있는 자는 관할 지방법원에 이의신청을 할 수 있다.
  \item 이의의 신청은 대법원규칙으로 정하는 바에 따라 등기소에 이의신청서를 제출하는 방법으로 한다.
  \item 관할 지방법원은 이의신청에 대하여 결정하기 전에 등기관에게 가등기 또는 이의가 있다는 뜻의 부기등기를 명령할 수 없다.
  \item 이의의 비용에 대하여는「비송사건절차법」을 준용한다.
  \item 이의에 대한 관할 지방법원의 결정에 대하여는「비송사건절차법」에 따라 항고할 수 있다.
\end{choices}
\end{question}
\answer{3}
